\documentclass{subfiles}
\begin{document}
    Todays cryptography is almost entirely based on public-key schemas.
        Broadly speaking, each of the entities involved in the communication 
        owns a public key \(e\), available by everyone, and a private key \(d\)
        that must be kept secret. To each key a transformation is associated,
        respectively to encrypt and decrypt the message.

    Here we focus on the most used public-key schemas, while in \Cref{sec:digital-signature}
        we describe how such schemas are used to garantee the autenticity of messages.
        Each of the schemas we are about to discuss is based on the hardness of some mathematical problem.
        Also, we denote by \(A\) and \(B\) two entities that want to communicate,
        and by \(E\) some unauthorized user.

    \subsection{RSA schema}
    \subfile{../sub/2.1 - RSA}
\end{document}
